\documentclass[a4paper]{article}
\pdfoutput=1
\usepackage[utf8]{inputenc}
\usepackage{jheppub}
\usepackage{overpic}
\usepackage{subcaption}
\usepackage{fancybox}
\usepackage{float}
\usepackage{amsfonts}
\usepackage{amsmath}
\usepackage{amsbsy}
\usepackage{graphicx}
\usepackage{amssymb}
\usepackage{amsthm}
\usepackage{bm}
\usepackage{comment}
\usepackage{epsfig}
\usepackage{latexsym}
\usepackage{pdflscape}
\usepackage{color}
\usepackage{slashed}
\usepackage{tikz}
\usetikzlibrary{arrows.meta,decorations.markings}

\tikzset{
  mid arrow/.style={
    postaction={
      decorate,
      decoration={
        markings,
        mark=at position 0.5 with {\arrow{Stealth[length=6pt,width=6pt]}}
      }
    }
  }
}

\usepackage{cancel}
\usepackage{xcolor}
\usepackage{braket}
\usepackage{empheq}
\usepackage{graphicx} % Allows including images
\usepackage{tikz-3dplot}
% réglage de la vue 3D : élévation=20°, azimut=110°
\tdplotsetmaincoords{20}{110}
\usetikzlibrary{decorations.pathreplacing,calc}
\usetikzlibrary{decorations.pathmorphing}

\definecolor{blue5}{rgb}{0, 0, 0.66666}
\definecolor{red4}{rgb}{0.8, 0, 0}
\definecolor{navyblue}{rgb}{0, 0.2745, 0.67843}
\definecolor{cornflowerblue}{rgb}{0.388235, 0.6941176, 0.898039}
\definecolor{forestgreen}{rgb}{0.13, 0.55, 0.13}
\usetikzlibrary{decorations.text,decorations.pathmorphing}

%Removes Header------------------------------------
\makeatletter
\def\@fpheader{\relax}
\makeatother



\usepackage{orcidlink}

%\DeclareMathOperator{\tr}{tr}
\DeclareMathOperator{\pf}{pf}
\DeclareMathOperator{\Ai}{Ai}
\DeclareMathOperator{\vol}{vol}
\DeclareMathOperator{\sgn}{sgn}
\DeclareMathOperator{\sn}{sn}
\DeclareMathOperator{\cn}{cn}
\DeclareMathOperator{\dn}{dn}
\DeclareMathOperator{\cd}{cd}
\DeclareMathOperator{\Det}{Det}
\DeclareMathOperator{\sech}{sech}
% \DeclareMathOperator{\tr}{tr}

\newcommand{\beq}{\begin{eqnarray}}
\newcommand{\eeq}{\end{eqnarray}}
\newcommand{\bea}{\begin{eqnarray}}
\newcommand{\eea}{\end{eqnarray}}
\newcommand{\be}{\begin{equation}}
\newcommand{\ee}{\end{equation}}
\newcommand{\bq}{\begin{equation}}
\newcommand{\eq}{\end{equation}}
\newcommand{\unit}{1\!\!1}
\newcommand{\half}{\frac{1}{2}}
\newcommand{\nn}{\nonumber}
\newcommand{\no}{\nonumber}
\def\tb{\tilde{\beta}}
\newcommand{\mpl}{M_{\rm Pl}}
\def\k{\kappa}

\newcommand{\bto}{\bar{\tau}^{\textbf{I}}}
\newcommand{\btt}{\bar{\tau}^{\textbf{II}}}




%%% Other definitions


\def\sp{\;\;\;,\;\;\;}
\def\l{\lambda}
%\def\z{\zeta}
\def\lab{\label}
\def\f{H}
\def\La{\Lambda}
\def\z{\zeta}
\def\t{\tau}
\def\e{\epsilon}
\def\o{\omega}
\def\le{\left}
\def\ri{\right}
\def\y{\psi}
\def\half{\frac12}
\def\q{\theta}
\def\cO{{\cal O}}
\def\m{\mu}
\def\n{\nu}
\def\td{\tilde}
\def\d{\delta}
\def\6{\partial}
\def\de{\partial}
\def\del{\nabla}
\def\ls{\ell_s}
\def\a{\alpha}
\def\b{\beta}
\def\lab{\label}
\def\parl{\parallel}
\def\tom{\tilde{\o}}
\def\dA{\dot{A}}
\def\dh{\dot{h}}
\def\df{\dot{f}}
\def\dX{\dot{X}}
\def\bo{\bar{\o}}
\def\r{x}
\def\bo{b_0}
\def\br{\bar{r}}
\def\dx{\delta X^\parl}
\def\dxy{\delta X^2}
\def\dxz{\delta X^3}
\def\ol{\overline}
\def\GM#1{\textit{[GM: #1]}}
\def\g{\gamma}
\def\s{\sigma}
\def\tr{{\rm Tr}}
\def\eps{\epsilon}
\def\fb{f}
\def\rd{\rm d}
\newcommand{\bom}[1]{\fboxsep2mm\fbox{
           $ \displaystyle{ #1} $}}
           
           \newcommand{\rar}{\rightarrow}
\newcommand{\bb}{\mathbb}
\newcommand{\mc}{\mathcal}
\newcommand{\im}{\mathrm{Im}}
\newcommand{\re}{\mathrm{Re}}
\newcommand{\mi}{\sqrt{|\mathcal{I}_4|}}


\def\lab{\label}
\def\le{\left}
\def\ri{\right}
\def\6{\partial}
\def\eps{\epsilon}


% Color definitions
\newcommand{\red}[1]{{\color{red} #1 \color{black}}}
\newcommand{\blue}[1]{{\color{blue} #1 \color{black}}}
\newcommand{\green}[1]{{\color{green} #1 \color{black}}}
\newcommand{\yellow}[1]{{\color{yellow} #1 \color{black}}}
\newcommand{\purple}[1]{{\color{purple} #1 \color{black}}}


% Comment Abbreviations
\newcommand{\IT}[1]{{\red{IG: #1}}} %Ioannis
\newcommand{\PB}[1]{{\blue{PB: #1}}} % Panos
\newcommand{\OP}[1]{{\green{OP: #1}}} % Olga
\newcommand{\PG}[1]{{\purple{PG: #1}}} % Paul

\title{Notes on boundary Liouville}

\author[a]{Panos Betzios\orcidlink{0000-0002-5350-9404},}
\author[b]{Paul Ghiringhelli,}
\author[b]{Olga Papadoulaki\orcidlink{0000-0001-5302-2930}}
\affiliation[a]{\href{https://www.ugent.be/we/physics-astronomy/en}{Department of Physics and Astronomy},
Ghent University, \\ Krijgslaan, 281-S9, 9000 Gent, Belgium}
\affiliation[b]{\href{https://www.cpht.polytechnique.fr/?q=en}{CPHT, CNRS, École polytechnique, Institut Polytechnique de Paris},  91120 Palaiseau, France}

\emailAdd{panos.betzios@ugent.be}
\emailAdd{paul.ghiringhelli@polytechnique.edu}
\emailAdd{olga.papadoulaki@polytechnique.edu}





\abstract{...}

\keywords{...}


\begin{document}
\maketitle
\flushbottom


%%%%%%%%%%%%%%%%%%%%%%%%%%%%%%%%%%%%%%%%%%%%%%%%%%%%%%%%%%

\section{Notations}
\paragraph{Central charge} We parametrize the central charge as
\begin{equation}
    c=1+6\,Q^2\,, \quad Q=b+b^{-1}\,.
\end{equation}
We do not assume anything particular range for the central charge. We want to obtain the shift equations as a functions of the momenta and the central charge.
\paragraph{Bulk conformal dimensions} We parametrize the conformal dimension of bulk primary fields $V_{P}(z,\bar{z})$ as 
\begin{equation}
  \Delta = \frac{Q^2}{4} - P^2\,, \quad P\in i \mathbb{R}_+\,.
\end{equation}

\paragraph{Boundary conformal dimension} We work on the upper half-plane. Boundary primary fields are denoted as $\Psi_{P}(x)^{\sigma_1\,\sigma_2}$ where $P$ is the momentum and $\sigma_{1,2}$ corresponds to the boundary conditions on the two sides of the puncture (in principle the two boundary cosmological constants.) We do not assume any particular rules for $\sigma_{i}$ (like FZZT boundary conditions).

\paragraph{Normalization} We work with normalizations of the two point functions/three point functions which are symmetric. Precisely, we write
\begin{equation}\label{eq:Normalizations}
\begin{aligned}
  \braket{V_{P_1}(0)V_{P_2}(1)}&=B(P)\, (\delta(P_1-P_2)+\delta(P_1+P_2))\,,\\
  \braket{V_{P_1}(0)V_{P_2}(1)V_{P_3}(\infty)}&= C(P_1,P_2,P_3)\,,\\
  \braket{\Psi_{P_1}^{\sigma_1 \sigma_2}(0)\Psi_{P_2}^{\sigma_2 \sigma_1}(1)}&= S_{P_1}^{\sigma_1 \sigma_2}\, (\delta(P_1-P_2)+\delta(P_1+P_2))\,,\\
  \braket{\Psi_{P_1}^{\sigma_1 \sigma_2}(0)\Psi_{P_2}^{\sigma_2 \sigma_3}(1)\Psi_{P_3}^{\sigma_3 \sigma_1}(\infty)} &= C_{P_1,P_2,P_3}^{\sigma_1 \sigma_2 \sigma_3}\,,\\
\end{aligned}
\end{equation}
and also 
\begin{equation}\label{eq:Normalizationsbis}
\begin{aligned}
  \braket{V_{P}(z,\bar{z})}&=\frac{U(P|\,\sigma)}{|z-\bar{z}|^{2\Delta}}\,.\\
  \braket{V_{P}(z,\bar{z})\Psi_{P'}^{\sigma \sigma}(x)}&=\frac{R(P;P'|\,\sigma)}{|z-\bar{z}|^{2\Delta -\Delta'}|z-x|^{2\Delta'}}
\end{aligned}
\end{equation}

\section{Equations from the Teschner trick}
We would like to write down all the possible equations one can get by insertion of appropriate degenerate fields. There are two things that we can do: either insert a bulk degenerate field $V_{\braket{2,1}}(z,\bar{z})$ which always decouple. This would lead to second order differential equations. We can also insert $\Psi_{\braket{3,1}}^{\sigma \sigma}(x)$ at the boundary. As was pointed out in the FZZT paper, $\Psi_{\braket{2,1}}^{\sigma \sigma}(x)$ decouples only for peculiar values of $\sigma$. The reason is that introducing a boundary modifies the ward identities obeyed by the boundary fields as the stress tensor at the boundary is perturbed by a certain primary operator which corresponds to the "line operator" (see FZZT paper: \PG{Add the expression of the boundary stress tensor in my convention}). In the other hand, there is a comprehensive reason for which $\Psi_{\braket{3,1}}^{\sigma \sigma}(x)$ always decouple. In fact, $V_{\braket{2,1}}(z,\bar{z})$ always decouple and as $\Im (z)\to 0$, the primary bulk operator can be expanded in boundary operators. The expansion should be compatible with the doubling trick\footnote{The origin of the doubling trick is encoded is the property of the stress tensor that satisfies $T(\bar{z})=\bar{T}(z)$} and thus the fusion rule of $V_{\braket{2,1}}(z)$ with itself (at $\bar{z}$ ) from which we obtain
\begin{equation}
  V_{\braket{2,1}}\,\times\,V_{\braket{2,1}} = \Psi_{\braket{1,1}}^{\sigma \sigma} + \Psi_{\braket{3,1}}^{\sigma \sigma}
\end{equation}

\paragraph{Three point function of boundary operators} We first consider the possibility of inserting a the bulk degenarate field $V_{\braket{2,1}}(z,\bar{z})$ in the boundary three point function with $\Re (z) \in(0,1)$
\begin{equation}
  \begin{aligned}
    &\braket{\Psi_{P_1}^{\sigma_1 \sigma_2}(0)V_{\braket{2,1}}(z,\bar{z})\Psi_{P_2}^{\sigma_2 \sigma_3}(1)\Psi_{P_3}^{\sigma_3 \sigma_1}(\infty)}\\
    &=\frac{R(P_{\braket{2,1}};P_{\braket{1,1}}|\sigma_2)}{S_{P_{\braket{1,1}}}^{\sigma_2 \sigma_2}}\,\frac{C_{P_1,P_{\braket{1,1}},P_1}^{\sigma_1 \sigma_2 \sigma_2}\,C_{P_1,P_2,P_3}^{\sigma_2 \sigma_2 \sigma_3}}{S_{P_1}^{\sigma_1 \sigma_2}}\,\mathcal{G}^{(1)}_{\braket{1,1}}(z,\bar{z})\\
    &+\frac{R(P_{\braket{2,1}};P_{\braket{3,1}}|\sigma_2)}{S_{P_{\braket{3,1}}}^{\sigma_2 \sigma_2}}\sum_{\eta\in\{{-1,0,1}\}}\frac{C_{P_1,P_{\braket{3,1}},P_1+\eta}^{\sigma_1 \sigma_2 \sigma_2}\,C_{P_1+\eta b,P_2,P_3}^{\sigma_1 \sigma_2 \sigma_3}}{S_{P_1+\eta b }^{\sigma_1 \sigma_2}}\,\mathcal{G}^{(1,\eta)}_{\braket{3,1}}(z,\bar{z})\\
    &=\frac{R(P_{\braket{2,1}};P_{\braket{1,1}}|\sigma_2)}{S_{P_{\braket{1,1}}}^{\sigma_2 \sigma_2}}\,\frac{C_{P_{\braket{1,1}},P_{2},P_2}^{\sigma_2 \sigma_2 \sigma_3}\,C_{P_1,P_2,P_3}^{\sigma_1 \sigma_2 \sigma_3}}{S_{P_2}^{\sigma_2 \sigma_3}}\,\mathcal{G}^{(2)}_{\braket{1,1}}(z,\bar{z})\\
    &+\frac{R(P_{\braket{2,1}};P_{\braket{3,1}}|\sigma_2)}{S_{P_{\braket{3,1}}}^{\sigma_2 \sigma_2}}\sum_{\eta\in\{{-1,0,1}\}}\frac{C_{P_{\braket{1,1}},P_{2},P_2+\eta b}^{\sigma_1 \sigma_2 \sigma_2}\,C_{P_2+\eta b,P_2,P_3}^{\sigma_1 \sigma_2 \sigma_3}}{S_{P_2+\eta b }^{\sigma_2 \sigma_3}}\,\mathcal{G}^{(2,\eta)}_{\braket{3,1}}(z,\bar{z})
  \end{aligned}
\end{equation}
The two equalities correspond to the different channel and this is the statement of crossing symmetry. The functions $\mathcal{G}$ correspond to the conformal blocks. I found the expansion \PG{To be checked!}
\begin{equation}
\begin{aligned}
  \mathcal{G}_{\braket{1,1}}^{(i)}(z,\bar{z})&=(z-\bar{z})^{\Delta'_{1,1}-2\Delta_{\braket{2,1}}}\,\left(\frac{z+\bar{z}}{2}\right)^{0} + ...\\
  \mathcal{G}_{\braket{3,1}}^{(i,\eta)}(z,\bar{z})&=(z-\bar{z})^{\Delta'_{3,1}-2\Delta_{\braket{2,1}}}\,\left(\frac{z+\bar{z}}{2}\right)^{\Delta'_{i+\eta}-\Delta'_{i}-\Delta'_{\braket{3,1}}} + ...
\end{aligned}
\end{equation}
We denoted by $\Delta'$ the conformal dimensions of the boundary primaries. The leading behavior is crucial to map them to the basis of the differential equation arising from the null vector decoupling. This generalized BPZ equation will lead to a second order (\PG{PDE I guess, since they effectively have 5 point by the doubling trick...}).

Alternatively, we can insert $\Psi^{\sigma_2 \sigma_2}_{\braket{3,1}}(x)$ with $x\in(0,1)$. One obtains
\begin{equation}
\begin{aligned}
    &\braket{\Psi_{P_1}^{\sigma_1 \sigma_2}(0)\Psi^{\sigma_2\sigma_2}_{\braket{3,1}}(x)\Psi_{P_2}^{\sigma_2 \sigma_3}(1)\Psi_{P_3}^{\sigma_3 \sigma_1}(\infty)}\\
    &=\sum_{\eta\in\{-1,0,1\}}\,\frac{C_{P_1,P_{\braket{3,1}},P_1+\eta b}^{\sigma_1 \sigma_2 \sigma_2}C_{P_1+\eta b,P_2,P_3}^{\sigma_1 \sigma_2 \sigma_3 }}{S_{P_1+\eta b}^{\sigma_1 \sigma_2}}\,\mathcal{G}^{(1,\eta)}(x)\\
    &=\sum_{\eta\in\{-1,0,1\}}\,\frac{C_{P_{\braket{3,1}},P_2,P_2+\eta b}^{\sigma_2 \sigma_2 \sigma_3}C_{P_1,P_2+\eta b,P_3}^{\sigma_1 \sigma_2 \sigma_3 }}{S_{P_2+ \eta b }^{\sigma_2 \sigma_3}}\,\mathcal{G}^{(2,\eta)}(x)
\end{aligned}
\end{equation}
The functions $\mathcal{G}$ are the conformal blocks. They have the following leading expansion
\begin{equation}
\begin{aligned}
  \mathcal{G}^{(1,\eta)}(x)\underset{x\to0}{\sim} &x^{\Delta'_{1+\eta}-\Delta'_{1}-\Delta'_{\braket{3,1}}}\\
  \mathcal{G}^{(2,\eta)}(x)\underset{x\to1}{\sim} &(1-x)^{\Delta'_{2+\eta}-\Delta'_{2}-\Delta'_{\braket{3,1}}}
\end{aligned}
\end{equation}
Unlike the previous case, the decoupling of $\Psi^{\sigma_2\sigma_2}_{\braket{3,1}}$ will produce a third order ODE (This is a ODE but we gain one order...). Let's emphasize this point: In bulk Liouville theory, the structure constants are such that they appropriate encode the connection formulae of a particular hypergeometric equation and also of higher order equations (arising from all the higher order degenerate fields decoupling). Boundary Liouville theory CFT data only encode the connection formulae of this third order equation (and those coming from fusing $\Psi^{\sigma_2\sigma_2}_{\braket{3,1}}$ with itself) while the decoupling of the null vector arise for specific values of the boundary cosmological constants. 

\paragraph{Two point function: 1 boundary operator with 1 bulk operator} It seems more convenient in this case to insert a bulk degenerate field.

\begin{equation}
  \begin{aligned}
    &\braket{V_{\braket{2,1}}(z,\bar{z})V_{P}(z_1,\bar{z}_1)\Psi_P'^{\sigma \sigma }(x)}\\
    &=\sum_{\epsilon = \pm} \frac{C(P,P_{2,1},P+\frac{\epsilon b}{2})}{B(P+\frac{\epsilon b}{2})}\,R(P+\frac{\epsilon b }{2};P'|\sigma) \mathcal{G}^{(\epsilon)}(z,z_1,x)\\
    &=\frac{R(P_{\braket{2,1}};P_{\braket{3,1}}|\sigma)}{S_{P_{\braket{3,1}}}^{\sigma \sigma}} \sum_{\eta \in \{-1,0,1\}}\frac{C_{P',P_{\braket{3,1},P'+\eta b}}^{\sigma \sigma \sigma }}{S_{P'+\eta b}^{\sigma \sigma}}R(P;P'+\eta b|\sigma)\mathcal{G}_{\braket{3,1}}^{(\eta)}(z,z_1,x)\\
    &+\frac{R(P_{\braket{2,1}};P_{\braket{1,1}}|\sigma)}{S_{P_{\braket{1,1}}}^{\sigma \sigma}} \frac{C_{P',P_{\braket{1,1},P'}}^{\sigma \sigma \sigma }}{S_{P'}^{\sigma \sigma}}R(P;P'|\sigma)\mathcal{G}_{\braket{1,1}}(z,z_1,x)\,.
  \end{aligned}
\end{equation}
\PG{I need to write the conformal blocks in terms of the invariants cross ratios (there should be two invariant cross ratios)}
The first equality corresponds to fusing the two bulk fields while the second corresponds to expanding the bulk field in terms of boundary operators.

\paragraph{1 point function}
We again insert a bulk degenerate field.
\begin{equation}
  \begin{aligned}
    \braket{V_{\braket{2,1}(z,\bar{z})}V_{P}(z,\bar{z})}
    &=\sum_{\epsilon = \pm} \frac{C_{P,P_{\braket{2,1}},P+\frac{\epsilon b }{2}}}{B(P+\frac{\epsilon b}{2})} U(P+\frac{\epsilon b}{2}|\sigma) \mathcal{G}^{(\epsilon)}(z,z_1)\\
    &=\frac{R(P_{\braket{2,1}};P_{\braket{3,1}}|\sigma)}{S_{P_{\braket{3,1}}}^{\sigma \sigma}} R(P;P_{\braket{3,1}}|\sigma)\mathcal{G}_{\braket{3,1}}(z,z_1)\\
    &+\frac{R(P_{\braket{2,1}};P_{\braket{1,1}}|\sigma)}{S_{P_{\braket{1,1}}}^{\sigma \sigma}}R(P;P_{\braket{1,1}}|\sigma)\mathcal{G}_{\braket{1,1}}(z,z_1)\,.
  \end{aligned}
\end{equation}
The first channel correspond to fusing the two bulk fields together while the second one correspond to expanding the bulk field in terms of boundary operators. 

\section{BPZ equation}
\paragraph{Bulk degenerate field}
The decoupling of the degenerate bulk field reads 
\begin{equation}
  \braket{\left(L_{-2}^{(z)}+\frac{1}{b^2}L_{-1}^{2\,(z)}\right)V_{\braket{2,1}}(z,\bar{z})\,X} = 0\,,
\end{equation}
where $X$ is a chain of fields inserted at different points (on the boundary or in the bulk).
\paragraph{Three point function of boundary operators}
We denote
\begin{equation}
  \mathcal{V}_{4}(z,\bar{z})=\braket{\Psi_{P_1}^{\sigma_1 \sigma_2}(0)V_{\braket{2,1}}(z,\bar{z})\Psi_{P_2}^{\sigma_2 \sigma_3}(1)\Psi_{P_3}^{\sigma_3 \sigma_1}(\infty)}\,.
\end{equation}
The constraint reads
\begin{equation}\label{eq:3pointfunctionBPZ1}
  0=\braket{\Psi_{P_1}^{\sigma_1 \sigma_2}(0)\left(L_{-2}^{(z)}+\frac{1}{b^2}L_{-1}^{2\,(z)}\right)V_{\braket{2,1}}(z,\bar{z})\Psi_{P_2}^{\sigma_2 \sigma_3}(1)\Psi_{P_3}^{\sigma_3 \sigma_1}(\infty)}\,.
\end{equation}
$L_{-1}$ just yields $\partial_z$ while $L_{-2}^{(z)}=\frac{1}{2\pi i}\int_{w \sim z}\frac{dw\,T(w)}{w-z}$. The standard trick is to deform the contour to pick up the residues at the other poles. In the meantime, we want to avoid factor of $\partial_{x_i}$ arising from the fact that 
\begin{equation}
  T(w)\Psi_{P}^{\sigma \sigma}(x_i)\sim \frac{\Delta_i'}{(w-x_{i})^2}+\frac{\partial_{x_{i}}}{w-x_{i}}\,.
\end{equation}
To do that we write
\begin{equation}
  \frac{1}{w-z}=\frac{w(w-1)}{z(z-1)}\frac{1}{w-z}-\frac{w-z}{z(z-1)}-\frac{2z-1}{z(z-1)}\,.
\end{equation}
This partial fraction decomposition has the benefit of removing one order of the pole at $w\sim0,w\sim1$ regarding the first term. The second term will correspond to the dilation operator $L_0$ while the last one will yield $L_{-1}=\partial_z$. Inserting this back to~\eqref{eq:3pointfunctionBPZ1} produces
\begin{equation}
\begin{aligned}
  0&=\frac{1}{b^2}\partial_z^2\,\mathcal{V}_{4}-\frac{\Delta_{2,1}}{z(z-1)}\mathcal{V}_{4}-\frac{2z-1}{z(z-1)}\partial_z \mathcal{V}_{4}\\
  &+\braket{\Psi_{P_1}^{\sigma 1 \sigma 2}(0)\frac{1}{2\pi i}\int_{w \sim z}\frac{w(w-1)}{z(z-1)}\frac{T(w)}{w-z} V_{\braket{2,1}}(z,\bar{z})\Psi_{P_2}^{\sigma 2 \sigma 3}(1)\Psi_{P_3}^{\sigma 3 \sigma 1}(\infty)}\,.
\end{aligned}
\end{equation}
We then deform the contour to pick up the poles at $0,1,\infty$ and $\bar{z}$. It yields
\begin{equation}
\begin{aligned}
  0&=\frac{1}{b^2}\partial_z^2\,\mathcal{V}_{4}-\frac{\Delta_{2,1}}{z(z-1)}\mathcal{V}_{4}-\frac{2z-1}{z(z-1)}\partial_z \mathcal{V}_{4}-\frac{\Delta_1'}{z^2(z-1)}\mathcal{V}_{4}+\frac{\Delta_2'}{z(1-z)^2}\mathcal{V}_{4}-\frac{\Delta'_3}{z(1-z)}\mathcal{V}_{4}\\
  &-\braket{\Psi_{P_1}^{\sigma 1 \sigma 2}(0)\frac{1}{2\pi i}\int_{w \sim \bar{z}}\frac{w(w-1)}{z(z-1)}\frac{T(w)}{w-z} V_{\braket{2,1}}(z,\bar{z})\Psi_{P_2}^{\sigma 2 \sigma 3}(1)\Psi_{P_3}^{\sigma 3 \sigma 1}(\infty)}\,.
\end{aligned}
\end{equation}
As 
\begin{equation}
\begin{aligned}
  T(w)V_{\braket{2,1}}(z,\bar{z}) &\sim \frac{\Delta_{2,1}V_{\braket{2,1}}(z,\bar{z})}{(w-z)^2}+\frac{\partial_zV_{\braket{2,1}}(z,\bar{z})}{w-z}+\frac{\Delta_{2,1}V_{\braket{2,1}}(z,\bar{z})}{(w-\bar{z})^2}+\frac{\partial_{\bar{z}}V_{\braket{2,1}}(z,\bar{z})}{w-\bar{z}}\,,\\
  \frac{w(w-1)}{z(z-1)(w-z)}&=\frac{\bar{z}(\bar{z}-1)}{z(z-1)}\frac{1}{\bar{z}-z}+\frac{(2\bar{z}-1)(\bar{z}-z)-\bar{z}(\bar{z}-1)}{z(z-1)(\bar{z}-z)^2}\times (w-z)+o(w-z)\,,
\end{aligned}
\end{equation}
I finally obtain after picking up the pole at $\bar{z}$, 
\begin{equation}
\begin{aligned}
  0&=\Bigg[\frac{1}{b^2}\partial_z^2-\frac{\Delta_{2,1}}{z(z-1)}-\frac{2z-1}{z(z-1)}\partial_z -\frac{\Delta_1'}{z^2(z-1)}+\frac{\Delta_2'}{z(1-z)^2}-\frac{\Delta'_3}{z(1-z)}\\
  &+\frac{\bar{z}(\bar{z}-1)\partial_{\bar{z}}}{z(z-1)(z-\bar{z})}-\frac{(2\bar{z}-1)(\bar{z}-z)-\bar{z}(\bar{z}-1)}{z(z-1)(\bar{z}-z)^2}\Delta_{\braket{2,1}}\Bigg]\mathcal{V}_{4}
\end{aligned}
\end{equation}
which can be rewritten as
\begin{equation}
\begin{aligned}
  0&=\Bigg[\frac{1}{b^2}\partial_z^2-\frac{2\Delta_{2,1}}{z(z-1)}-\frac{2z-1}{z(z-1)}\partial_z -\frac{\Delta_1'}{z^2(z-1)}+\frac{\Delta_2'}{z(1-z)^2}-\frac{\Delta'_3}{z(1-z)}\\
  &+\frac{\bar{z}(\bar{z}-1)\partial_{\bar{z}}}{z(z-1)(z-\bar{z})}-\frac{\Delta_{2,1}}{(z-\bar{z})^2}\Bigg]\mathcal{V}_{4}
\end{aligned}
\end{equation}
As expected, it yields a second order PDE. This equation should also be obeyed by all the conformal blocks cited above.

\end{document}

